%----------------------------------------------------------------------
\section{Self-Enrichment and Topos Structure}\label{sec:topos}
%----------------------------------------------------------------------

The monoidal closed structure of Section \ref{sec:monoidal}
equips $\mathcal{C}(S,\kappa)$ with internal hom-objects.
We now show that the category is canonically enriched over
itself, that the three ingredients of enrichment correspond
to the three vertices of the $\GF(2)$ toroid, and that the
resulting structure satisfies the axioms of a topos with a
four-valued internal logic.

\subsection{Self-enrichment from triadic closure}

\begin{definition}[Enrichment data]\label{def:enrichment-data}
An enrichment of a category $\mathcal{C}$ over a monoidal
category $(\mathcal{V}, \otimes, I)$ requires:
\begin{enumerate}[label=(\roman*),nosep]
  \item For each pair of objects $A, B$, a \emph{hom-object}
    $[A,B] \in \mathcal{V}$.
  \item For each triple $A, B, C$, a \emph{composition
    morphism} $\circ_{A,B,C}: [B,C] \otimes [A,B] \to [A,C]$
    in $\mathcal{V}$.
  \item For each object $A$, an \emph{identity morphism}
    $j_A: I \to [A,A]$ in $\mathcal{V}$.
\end{enumerate}
These must satisfy associativity and unitality coherence
conditions.
\end{definition}

\begin{theorem}[Canonical self-enrichment]\label{thm:self-enrich}
$\mathcal{C}(S,\kappa)$ is canonically enriched over
$(\mathcal{C}(S,\kappa), \otimes, I)$, with the three
ingredients of enrichment corresponding to the three
vertices of the $\GF(2)$ toroid:
\begin{itemize}[nosep]
  \item \textbf{Objects} (the things being enriched): the
    $\kappa$-perspective. Objects are $\kappa$-equivalence
    classes.
  \item \textbf{Hom-objects} $[A,B]$: the $\sigma$-perspective.
    Each $[A,B]$ is a sub-collection of $\kappa$ (the set
    of discriminators witnessing $A \to B$), viewed as an
    object via the $\sigma$-regime.
  \item \textbf{Composition morphisms}: the $\tau$-perspective.
    Composition $\circ_{A,B,C}: [B,C] \otimes [A,B] \to [A,C]$
    is a $\tau$-discriminator that acts on pairs of
    compatible discriminators (one from $[B,C]$, one from
    $[A,B]$) and produces their sequential composite
    in $[A,C]$.
\end{itemize}
\end{theorem}

\begin{proof}
\emph{Hom-objects.}
$[A,B]$ is the set $\{d \in \kappa : \tau_d(A) = 1,\;
\sigma_d(B) = 1\}$ (Definition \ref{def:internal-hom}).
Under the $\sigma$-perspective (Proposition
\ref{prop:perspectives}), this set is a sub-population of
$\kappa$, hence a legitimate object of the category when
viewed from the $\sigma$-vertex.

\emph{Composition.}
Given $d_1 \in [A,B]$ ($\tau_{d_1}(A) = 1$,
$\sigma_{d_1}(B) = 1$) and $d_2 \in [B,C]$
($\tau_{d_2}(B) = 1$, $\sigma_{d_2}(C) = 1$), the
sequential filtration through $d_1$ then $d_2$ produces
a composite discrimination placing $A$ in $\tau$ and $C$
in $\sigma$ (Definition \ref{def:composition}). The map
$(d_2, d_1) \mapsto d_2 \circ d_1 \in [A,C]$ is a
morphism in the $\sigma$-perspective: it is a
$\tau$-discriminator that separates ``compatible pairs of
discriminators that compose to a given target'' from
``incompatible pairs.''

\emph{Identity.}
The identity morphism $j_A: I \to [A,A]$ sends the
monoidal unit $I$ (the degenerate grade-0 element) to
the identity discriminator $\mathrm{id}_A = d \wedge d = 0
\in [A,A]$. This is a morphism from the trivial object
to the self-hom, selecting the degenerate self-discrimination.

\emph{Associativity.}
Composition is associative because sequential filtration is:
$(d_3 \circ d_2) \circ d_1 = d_3 \circ (d_2 \circ d_1)$,
inherited from associativity of the wedge product
(Theorem \ref{thm:cat-axioms}).

\emph{Unitality.}
$\mathrm{id}_B \circ d_1 = d_1$ and $d_2 \circ \mathrm{id}_B
= d_2$, since composing with the degenerate discriminator
(grade 0) leaves any discriminator unchanged.

All enrichment coherences are satisfied. The enrichment is
\emph{canonical}: no choices are made. The three ingredients
are the three vertices of the triangle identity, and the
coherences are properties of the exterior algebra over
$\GF(2)$.
\end{proof}

\begin{remark}[Self-enrichment is triadic closure]
The self-enrichment of $\mathcal{C}(S,\kappa)$ is not an
additional property discovered about the category. It is
the triadic closure of the $\GF(2)$ toroid: the fact that
all three vertices are equally valid perspectives on the
same structure means the category can describe its own
morphism spaces using its own objects and morphisms. A
category that is \emph{not} self-enriched would require one
vertex of the triangle to be privileged or external, which
the toroidal symmetry forbids.
\end{remark}

\subsection{The subobject classifier}\label{sec:subobject}

\begin{definition}[Subobject]\label{def:subobject}
A \emph{subobject} of an object $A$ in
$\mathcal{C}(S,\kappa)$ is a sub-collection of $A$
obtained by applying a discriminator $d$ and selecting
a $\tau/\sigma$ class. The subobject is specified by
the $\GF(2) \times \GF(2)$ profile of $d$ restricted
to $A$.
\end{definition}

\begin{definition}[Subobject classifier]\label{def:subobj-class}
The \emph{subobject classifier} of $\mathcal{C}(S,\kappa)$
is the object
\begin{equation}
  \Omega = \GF(2) \times \GF(2) = \{(0,0),\; (1,0),\;
  (0,1),\; (1,1)\},
\end{equation}
the four-cell structure of Definition \ref{def:discriminator}.
For any subobject $m: S' \hookrightarrow A$, there exists a
unique \emph{classifying morphism} $\chi_m: A \to \Omega$
that sends each element $a \in A$ to its $(\tau_d(a),
\sigma_d(a))$ profile under the discriminator $d$ that
defines the subobject.
\end{definition}

\begin{theorem}[$\Omega$ is a subobject classifier]
\label{thm:subobj}
$\Omega = \GF(2) \times \GF(2)$ satisfies the universal
property of a subobject classifier: for every monomorphism
$m: S' \hookrightarrow A$, there exists a unique
$\chi_m: A \to \Omega$ such that $S'$ is the pullback of
the ``true'' morphism $\top: 1 \to \Omega$ along $\chi_m$.
\end{theorem}

\begin{proof}
The ``true'' morphism selects the value $(1,0) \in \Omega$
(classified under $\tau$, not classified under $\sigma$).
Given a subobject $S' \hookrightarrow A$ defined by a
discriminator $d$ with $S' = \{a \in A : \tau_d(a) = 1,\;
\sigma_d(a) = 0\}$, the classifying morphism
$\chi_m: A \to \Omega$ sends each $a$ to
$(\tau_d(a), \sigma_d(a))$.

\emph{Existence}: $\chi_m$ is defined by the discriminator
$d$ that determines $S'$. Each element of $A$ receives a
unique $(\tau, \sigma)$ profile.

\emph{Uniqueness}: any morphism $\chi: A \to \Omega$ that
pulls back $\top$ to $S'$ must agree with $\chi_m$ on
$S'$ (sending elements to $(1,0)$) and must send elements
of $A \setminus S'$ to values other than $(1,0)$. The
specific values for $A \setminus S'$ are determined by
$d$'s $\tau/\sigma$ profile on those elements, which is
unique for a given discriminator.

\emph{Pullback}: $S' = \chi_m^{-1}(\{(1,0)\}) = \chi_m^{-1}(\top)$,
which is the pullback of $\top$ along $\chi_m$.
\end{proof}

\begin{remark}[Four-valued internal logic]
The subobject classifier $\Omega = \GF(2) \times \GF(2)$
has four elements, so the internal logic of
$\mathcal{C}(S,\kappa)$ is \emph{four-valued}:
\begin{center}
\begin{tabular}{c|l|l}
Value & $(\tau, \sigma)$ & Logical reading \\
\hline
$(1,0)$ & in $\tau$, not in $\sigma$ & True \\
$(0,1)$ & not in $\tau$, in $\sigma$ & False \\
$(0,0)$ & neither & Unknown (gap) \\
$(1,1)$ & both & Overdetermined (overlap)
\end{tabular}
\end{center}
This is Belnap's four-valued logic, arising constructively
from the discrimination framework. The four values are not
imposed by fiat; they are the four cells of the
$\GF(2) \times \GF(2)$ membership profile that has been
present since Definition \ref{def:discriminator}.

Boolean two-valued logic is recovered when the logic
is restricted to the subset $\{(1,0), (0,1)\}$ --- the
Boolean dimensions of Definition \ref{def:boolean-dim},
where every element is assigned to exactly one of $\tau$
or $\sigma$. As established throughout this paper,
Boolean structure is a local achievement, not a global
given. The internal logic reflects this: it defaults to
four values and specializes to two where Booleanness has
been earned.
\end{remark}

\subsection{Topos structure}

The emergent category $\mathcal{C}(S,\kappa)$ for a
\emph{fixed} $\kappa$ has a subobject classifier and power
objects, but may lack some finite limits (if $\kappa$ does
not contain enough discriminators to compute them). Rather
than claiming topos structure for a fixed $\kappa$ and
patching with a pro-system, we construct the topos
directly as a fibered category over the poset of
$\kappa$-states.

\begin{definition}[The fibered category]\label{def:fibered}
Let $\mathbf{K}$ be the category whose objects are
$\kappa$-states (ordered by sub-regime inclusion
$\kappa' \subseteq \kappa$) and whose morphisms are the
inclusions. The assignment
$\kappa \mapsto \mathcal{C}(S,\kappa)$ is a pseudofunctor
$\mathbf{K} \to \mathbf{Cat}$ (covariant via the free
functors $F: \mathcal{C}(S,\kappa') \to
\mathcal{C}(S,\kappa)$).

The \emph{fibered category} $\int \mathcal{C}$ is the
Grothendieck construction on this pseudofunctor. An object
of $\int \mathcal{C}$ is a pair $(\kappa, A)$ where $A$
is an object of $\mathcal{C}(S,\kappa)$ --- an equivalence
class together with the regime at which it is defined. A
morphism $(\kappa, A) \to (\kappa', B)$ is an inclusion
$\kappa \subseteq \kappa'$ together with a morphism
$F(A) \to B$ in the finer category.
\end{definition}

\begin{theorem}[Fibered topos]\label{thm:topos}
The fibered category $\int \mathcal{C}$ is an elementary
topos.
\end{theorem}

\begin{proof}
\emph{Finite limits.}
Let $D$ be a finite diagram in $\int \mathcal{C}$ with
components at stages $\kappa_1, \ldots, \kappa_m$. Define
$\kappa_\vee = \kappa_1 \vee \cdots \vee \kappa_m$: the
smallest regime containing all discriminators from all
$\kappa_i$. Since $D$ is finite, $\kappa_\vee$ is finite.

At the fiber over $\kappa_\vee$, all components of $D$
are present (embedded via the free functors). Proposition
\ref{prop:limits} establishes that limits exist in
$\mathcal{C}(S,\kappa_\vee)$ whenever $\kappa_\vee$
contains sufficient discriminators. If the limit does not
exist at $\kappa_\vee$, the obstruction is a residue
(Proposition \ref{prop:limits}), which triggers a
$\kappa$-evolution to $\kappa_\vee'$ where the limit does
exist. This evolution is finite (finitely many residues
for a finite diagram). The limit of $D$ in
$\int \mathcal{C}$ is the pair $(\kappa_\vee', L)$ where
$L$ is the limit computed at $\kappa_\vee'$.

Products: the product of $(\kappa_1, A)$ and
$(\kappa_2, B)$ is $(\kappa_1 \vee \kappa_2,\; A \times B)$,
where $A \times B$ is the joint equivalence class at the
finer regime. Equalizers: computed at the join, as the
sub-class on which two morphisms agree. Terminal object:
$(\kappa_\varnothing, \top)$ where $\kappa_\varnothing$
is the empty regime (zero dimensions, one equivalence
class containing everything).

\emph{Subobject classifier.}
$\Omega = \GF(2) \times \GF(2)$ does not depend on
$\kappa$ (the four-cell structure is the same at every
resolution). Therefore $\Omega$ defines a \emph{constant
object} in $\int \mathcal{C}$: the family
$\{(\kappa, \Omega)\}_{\kappa \in \mathbf{K}}$ with
identity transition maps. Theorem \ref{thm:subobj}
establishes the universal property at each fiber, and
the constant structure ensures coherence across fibers.

\emph{Power objects.}
We verify the universal property: for every relation
$R \hookrightarrow A \times B$ (a monomorphism in the
fiber over $\kappa$), there exists a unique
$\chi_R: A \to \Omega^B$ such that $R$ is the pullback
of the membership relation
$\in_B: B \times \Omega^B \to \Omega$ along
$\mathrm{id}_B \times \chi_R$.

The power object is $\mathcal{P}(B) = \Omega^B$: the
object of morphisms $B \to \Omega$ in
$\mathcal{C}(S,\kappa)$. An element of $\Omega^B$ is a
function $\phi: B \to \GF(2) \times \GF(2)$ assigning
each $b \in B$ a four-valued membership profile. Note
that $\Omega$ is four-valued, not two-valued, so
$\Omega^B$ is larger than the classical power set: it
includes ``subobjects'' with gaps and overlaps.

Given a relation $R \hookrightarrow A \times B$, define
$\chi_R: A \to \Omega^B$ by
$\chi_R(a)(b) = (\tau_R(a,b), \sigma_R(a,b))$, where
$\tau_R(a,b) = 1$ iff $(a,b)$ is in the $\tau$-class
of $R$'s defining discriminator, and similarly for
$\sigma_R$.

\emph{Existence:} $\chi_R$ is well-defined because $R$'s
discriminator produces a $(\tau, \sigma)$ profile for each
$(a,b)$ pair. For fixed $a$, the map $b \mapsto
\chi_R(a)(b)$ is an element of $\Omega^B$.

\emph{Pullback:} The fiber of $\chi_R$ over the ``true''
element $\top = (1,0) \in \Omega$ is
$\{(a,b) : \chi_R(a)(b) = (1,0)\} = \{(a,b) :
\tau_R(a,b) = 1, \sigma_R(a,b) = 0\}$, which is the
Boolean part of $R$. Over $\Omega$, the pullback recovers
the full four-valued relation $R$.

\emph{Uniqueness:} Any $\chi': A \to \Omega^B$ that pulls
back $\in_B$ to $R$ must satisfy
$\chi'(a)(b) = (\tau_R(a,b), \sigma_R(a,b))$ for all
$(a,b)$ (since the pullback condition determines each
component). Therefore $\chi' = \chi_R$.

\emph{Representability:} $\Omega^B$ must be an object of
the category $\mathcal{C}(S,\kappa)$, not merely a
set-theoretic construction. By monoidal closure
(Theorem \ref{thm:closure}), the internal hom $[B, \Omega]$
exists as an object of $\mathcal{C}(S,\kappa)$: it is
the collection of discriminators $d \in \kappa$ such that
$d$'s $\tau/\sigma$ profile restricted to $B$ takes values
in $\Omega$ (Definition \ref{def:internal-hom}). This
internal hom is $\Omega^B$. The representability is a
consequence of the closed monoidal structure, not an
additional assumption.

The free functors $F: \mathcal{C}(S,\kappa') \to
\mathcal{C}(S,\kappa)$ preserve $\Omega$ (it is constant
across fibers) and preserve finite limits (Proposition
\ref{prop:forgetful}), so they preserve the exponential
$\Omega^B$. Power objects are well-defined in the fibered
topos $\int \mathcal{C}$.
\end{proof}

\begin{remark}[Evolution as structure, not narrative]
The fibered construction makes $\kappa$-evolution into
first-class categorical structure rather than an external
narrative about a changing topos. Objects in
$\int \mathcal{C}$ carry their resolution: $(\kappa, A)$
records not just \emph{what} $A$ is but \emph{at what
level of discrimination} it was identified. Morphisms
between different fibers are the free functors --- the
embeddings from coarser to finer regimes that the
framework has used throughout.

Finite limits are computed at finite joins. No appeal to
a completed colimit or an ``infinite $\kappa$'' is needed
for any finite-limit claim. The operationalist axiom
(Definition \ref{def:operationalist}) is invoked only for
claims about infinite limits and completeness, which is the
honest boundary.
\end{remark}

\begin{corollary}[Pro-system interpretation]
\label{cor:pro-topos}
The fibered topos $\int \mathcal{C}$ can equivalently be
viewed as a pro-system (filtered diagram) of topoi
$\{\mathcal{C}(S,\kappa_n)\}$ connected by geometric
morphisms. As $\kappa$ evolves through states
$\kappa_1 \subset \kappa_2 \subset \cdots$, the free
functors (left exact left adjoints) provide geometric
morphisms. The filtered colimit $\varinjlim
\mathcal{C}(S,\kappa_n)$ is $\kappa$-indistinguishable
from the pro-system (Proposition \ref{prop:infinity}).
\end{corollary}

\begin{remark}[A non-Boolean topos]
$\mathcal{C}(S,\kappa)$ is a topos whose internal logic is
not Boolean but four-valued. This places it in the family
of non-Boolean topoi studied in constructive mathematics
and sheaf theory, but with a distinctive character: the
four-valuedness arises not from open-set lattices or
forcing conditions but from the $\GF(2) \times \GF(2)$
structure of discrimination itself.

The logical connectives of the internal logic are
determined by operations on $\Omega$:
\begin{itemize}[nosep]
  \item \textbf{Conjunction} $\wedge$: componentwise
    multiplication in $\GF(2) \times \GF(2)$.
  \item \textbf{Disjunction} $\vee$: componentwise
    maximum (where $1 > 0$).
  \item \textbf{Negation} $\neg$: swap the two components,
    $\neg(\tau,\sigma) = (\sigma, \tau)$. Note that
    $\neg\neg(\tau,\sigma) = (\tau,\sigma)$: double negation
    elimination holds. But $(\tau, \sigma) \vee
    \neg(\tau,\sigma) = (\tau,\sigma) \vee (\sigma,\tau)$,
    which equals $(\max(\tau,\sigma), \max(\sigma,\tau))$.
    This equals $(1,1)$ (overlap) only when both $\tau$
    and $\sigma$ are non-zero. For the gap $(0,0)$,
    $p \vee \neg p = (0,0)$: excluded middle \emph{fails}.
  \item \textbf{Implication} $\Rightarrow$: the internal
    hom in $\Omega$, determined by the closure adjunction
    (Theorem \ref{thm:closure}).
\end{itemize}
Excluded middle fails at the gap but holds everywhere
else. This is the logical expression of the framework's
central thesis: Boolean structure is a local achievement.
The gap is where Booleanness has not been earned, and the
internal logic faithfully records this.
\end{remark}

\begin{remark}[The topos as the natural home]
The topos $\mathcal{C}(S,\kappa)$ is the natural universe
of discourse for the discrimination framework. It provides:
\begin{itemize}[nosep]
  \item An internal logic (four-valued, locally Boolean)
    that matches the framework's epistemology.
  \item Self-enrichment (Theorem \ref{thm:self-enrich})
    that allows the category to describe its own structure.
  \item Monoidal closure (Theorem \ref{thm:closure}) that
    connects the $n$-categorical and monoidal perspectives
    via the toroid.
  \item A subobject classifier that \emph{is} the
    four-cell structure from which the entire framework
    was built.
\end{itemize}
The topos is not an additional layer placed on top of the
discrimination framework. It is the discrimination
framework, viewed categorically. The four-cell structure,
which began as the membership profile of a single
discriminator (Definition \ref{def:discriminator}), turns
out to be the subobject classifier of the entire category
--- the logical atom from which the topos's internal
universe is generated.
\end{remark}

\subsection{Proof theory of the internal logic}\label{sec:proof-theory}

We now characterize the internal logic of
$\mathcal{C}(S,\kappa)$ precisely, clarifying its
relationship to Belnap's four-valued logic and its
behavior under $\kappa$-evolution.

\begin{definition}[Evidence semantics]\label{def:evidence}
The two components of the membership profile
$(\tau_d(a), \sigma_d(a)) \in \GF(2) \times \GF(2)$
are \emph{independent partial assertions}:
\begin{itemize}[nosep]
  \item $\tau_d(a) = 1$: positive evidence that predicate
    $\tau_d$ fires on $a$. A positive assertion of
    $\tau$-membership.
  \item $\tau_d(a) = 0$: the predicate $\tau_d$ does not
    fire on $a$. This is a positive fact about the
    predicate's behavior, \emph{not} a negative assertion
    about $a$'s membership. Absence of evidence is not
    evidence of absence.
  \item $\sigma_d(a) = 1$ and $\sigma_d(a) = 0$:
    analogously for $\sigma$.
\end{itemize}
The value $0$ is \emph{silence}, not falsehood. A
silent predicate makes no assertion about the element;
it reports only that this predicate, applied to this
element, produced no match.
\end{definition}

\begin{definition}[Information ordering]\label{def:info-order}
The four cells carry a natural \emph{information ordering}
$\sqsubseteq$ (the Belnap lattice):
\begin{equation}
  (0,0) \sqsubseteq (1,0), \qquad
  (0,0) \sqsubseteq (0,1), \qquad
  (1,0) \sqsubseteq (1,1), \qquad
  (0,1) \sqsubseteq (1,1).
\end{equation}
The bottom element $(0,0)$ is maximally uninformative:
both predicates are silent. The top element $(1,1)$ is
maximally informative but conflicted: both predicates
fire. The singly informative states $(1,0)$ and $(0,1)$
are the stable, Boolean-compatible middle.

The ordering is:
\begin{center}
$(0,0)$ --- gap (no evidence) \\[2pt]
$\swarrow \qquad \searrow$ \\[2pt]
$(1,0)$ \qquad\quad $(0,1)$ --- Boolean states \\[2pt]
$\searrow \qquad \swarrow$ \\[2pt]
$(1,1)$ --- overlap (conflicting evidence)
\end{center}
\end{definition}

\begin{proposition}[Residues are extremal]\label{prop:extremal}
The two residue types (Definition \ref{def:residue},
Remark \ref{rem:residue-origins}) correspond to the two
extremes of the information ordering:
\begin{itemize}[nosep]
  \item Gap residues $(0,0)$: the information-minimal state.
    Neither predicate fires. The element is invisible to
    the discriminator.
  \item Overlap residues $(1,1)$: the information-maximal
    state. Both predicates fire. The element is above the
    discriminator's resolution.
\end{itemize}
Both have XOR signature $\chi_d(a) = 0$ in $\GF(2)$. Both
are unstable: they trigger $\kappa$-evolution. The stable
states $(1,0)$ and $(0,1)$ have $\chi_d(a) = 1$ and do not
trigger evolution.
\end{proposition}

\begin{theorem}[Dynamic proof theory]\label{thm:proof-theory}
The proof theory of the internal logic has two layers:
a static layer (valid within a fixed $\kappa$-state) and
a dynamic layer (governing transitions between
$\kappa$-states).

\emph{Static layer.}
Within a fixed $\kappa$, the logic is Belnap's
First Degree Entailment (FDE). The valid inference rules
are:
\begin{enumerate}[label=(\roman*),nosep]
  \item \textbf{Conjunction introduction}: from
    $(\tau_1, \sigma_1)$ and $(\tau_2, \sigma_2)$, infer
    $(\tau_1 \cdot \tau_2,\; \sigma_1 \cdot \sigma_2)$.
    (Both predicates must independently fire for the
    conjunction to fire.)
  \item \textbf{Disjunction introduction}: from
    $(\tau_1, \sigma_1)$, infer
    $(\tau_1, \sigma_1) \vee (\tau_2, \sigma_2) =
    (\max(\tau_1, \tau_2),\; \max(\sigma_1, \sigma_2))$.
    (Either predicate firing suffices.)
  \item \textbf{Modus ponens}: from $(\tau, \sigma) = (1,0)$
    and the discriminator encoding the implication, infer
    the consequent. Modus ponens requires the premise to be
    in the Boolean state $(1,0)$; it does not fire from
    $(0,0)$ (silence) or $(1,1)$ (conflict).
  \item \textbf{Modus tollens does not hold in general.}
    From $\sigma_d(b) = 0$ (the consequent's $\sigma$-predicate
    is silent) and an implication $a \to b$, one cannot
    infer $\tau_d(a) = 0$. Silence in the consequent means
    no evidence was found, not that the consequent is false.
    Modus tollens requires Boolean structure
    (Definition \ref{def:boolean-dim}) on the relevant
    dimension.
\end{enumerate}

\emph{Dynamic layer.}
When an inference produces or encounters a residue
(a $(0,0)$ or $(1,1)$ state), the static logic provides
no further deductions. Instead, the residue triggers
$\kappa$-evolution:
\begin{enumerate}[label=(\roman*),nosep]
  \item The residue specifies the discriminator needed
    to resolve it (Theorem \ref{thm:exhaustive}).
  \item $\kappa$ evolves by consuming the residue.
  \item In the new $\kappa$-state, the formerly extremal
    element may now occupy a Boolean state $(1,0)$ or
    $(0,1)$, enabling further static deductions.
\end{enumerate}
The dynamic layer drives the logic toward Boolean states:
each evolution resolves a residue, moving an element from
the extremes toward the informative middle.
\end{theorem}

\begin{proof}
\emph{Static layer.}
The connectives defined on $\Omega = \GF(2) \times \GF(2)$
in the non-Boolean topos remark (Section \ref{sec:topos})
match exactly the connectives of Belnap's FDE. The truth
tables are identical:
conjunction is componentwise multiplication,
disjunction is componentwise maximum,
and negation swaps components. FDE is the logic of the
Belnap lattice, and $\GF(2) \times \GF(2)$ \emph{is} the
Belnap lattice.

Modus ponens holds for Boolean premises because $(1,0)$
provides unambiguous positive evidence in $\tau$ and
unambiguous silence in $\sigma$, so the implication
discriminator can chain. Modus tollens fails for
non-Boolean states because silence ($0$) in the consequent
is not negation --- it carries no information from which
the premise's status can be deduced.

\emph{Dynamic layer.}
By Proposition \ref{prop:extremal}, residues occupy the
extremes of the information ordering. By Theorem
\ref{thm:exhaustive}, each residue specifies its own
resolution. By Proposition \ref{prop:monotone},
$\kappa$-evolution is monotonic: resolved states are
never reverted. Therefore the dynamic layer moves
elements along the information ordering from extremes
toward the Boolean middle.
\end{proof}

\begin{remark}[Convergence toward Booleanness]
The dynamic proof theory describes a logic that
\emph{tends toward Boolean} without ever being required
to arrive. Each $\kappa$-evolution resolves one or more
residues, converting $(0,0)$ or $(1,1)$ states into
$(1,0)$ or $(0,1)$ states on new dimensions. But new
discriminations can introduce new residues on their own
dimensions. The logic oscillates:
\begin{itemize}[nosep]
  \item Resolution: consuming a residue earns Boolean
    structure on one dimension.
  \item Introduction: articulating a new predicate may
    produce new residues on new dimensions.
\end{itemize}
The net effect is that the proportion of Boolean dimensions
grows as $\kappa$ evolves (resolved dimensions stay
resolved), but full Booleanness across all dimensions
is an asymptotic condition, not a terminus.
\end{remark}

\begin{remark}[Relationship to classical, intuitionistic,
  and paraconsistent logics]
The internal logic of $\mathcal{C}(S,\kappa)$ relates to
the standard logical traditions as follows:
\begin{itemize}[nosep]
  \item \textbf{Classical logic} is the restriction to
    Boolean dimensions (Definition \ref{def:boolean-dim}),
    where $\Omega$ collapses to $\{(1,0), (0,1)\}$.
    Excluded middle holds. Modus tollens holds. This is a
    local achievement, valid on specific dimensions.
  \item \textbf{Intuitionistic logic} shares the failure of
    excluded middle but for a different reason.
    Intuitionistically, $p \vee \neg p$ fails because
    disjunction requires a constructive witness.
    In the discrimination logic, $p \vee \neg p$ fails at
    the gap because \emph{neither $p$ nor $\neg p$ has
    evidence}. The failure is epistemic (no evidence),
    not constructive (no witness).
  \item \textbf{Paraconsistent logic} tolerates
    contradictions. The overlap state $(1,1)$ is a
    ``contradiction'' --- both $\tau$ and $\sigma$ fire
    --- but it does not produce explosion (derivation of
    arbitrary conclusions). In FDE, $p \wedge \neg p$ does
    not entail $q$. The logic is paraconsistent at the
    overlap.
\end{itemize}
The discrimination logic therefore unifies aspects of
all three traditions: it is locally classical (on Boolean
dimensions), shares intuitionistic features at the gap,
and is paraconsistent at the overlap. The unification
arises not from a philosophical commitment but from the
$\GF(2) \times \GF(2)$ structure of the subobject
classifier.
\end{remark}

\begin{proposition}[First-order logic is internal]
\label{prop:fol}
First-order logic (FOL) --- including its quantifiers,
inference rules, and proof theory --- is derivable as a
fragment of the internal logic of the topos
$\mathcal{C}(S,\kappa)$. The metalanguage required to
read and validate the framework is not external to it.
\end{proposition}

\begin{proof}
The internal logic of an elementary topos is a
higher-order intuitionistic type theory. FOL is a
fragment. We derive each component:

\emph{Propositional connectives.}
$\wedge$, $\vee$, $\neg$, and $\Rightarrow$ are
operations on $\Omega = \GF(2) \times \GF(2)$, already
defined (Theorem \ref{thm:proof-theory}).

\emph{Quantifiers.}
Let $f: A \to B$ be a morphism in $\mathcal{C}(S,\kappa)$
and let $P: A \to \Omega$ be a predicate on $A$ (a
subobject classifier morphism).
\begin{itemize}[nosep]
  \item The \emph{universal quantifier}
    $\forall_f: \Omega^A \to \Omega^B$ is the right adjoint
    to the pullback functor $f^*: \Omega^B \to \Omega^A$.
    $(\forall_f P)(b) = \bigwedge_{a \in f^{-1}(b)} P(a)$:
    the conjunction (componentwise multiplication in
    $\GF(2) \times \GF(2)$) of $P$'s values over the
    fiber of $f$ above $b$. This is a morphism in the
    topos, not an external operation.
  \item The \emph{existential quantifier}
    $\exists_f: \Omega^A \to \Omega^B$ is the left adjoint
    to $f^*$.
    $(\exists_f P)(b) = \bigvee_{a \in f^{-1}(b)} P(a)$:
    the disjunction (componentwise maximum) over the fiber.
\end{itemize}

\emph{Inference rules.}
\begin{itemize}[nosep]
  \item \textbf{Modus ponens} is the evaluation morphism
    $\mathrm{ev}: [A, B] \otimes A \to B$
    (Definition \ref{def:eval}), applied to the case
    $A = \Omega$, $B = \Omega$. Given an implication
    $p \Rightarrow q$ (an element of $[\Omega, \Omega]$)
    and a premise $p$ (an element of $\Omega$), evaluation
    produces the conclusion $q$. This is a morphism in the
    topos.
  \item \textbf{Universal instantiation} is the counit of
    the $\forall_f \dashv f^*$ adjunction: given
    $\forall_f P$ and an element $a \in f^{-1}(b)$, the
    counit produces $P(a)$.
  \item \textbf{Existential generalization} is the unit of
    the $f^* \dashv \exists_f$ adjunction: given $P(a)$,
    the unit produces $\exists_f P(f(a))$.
\end{itemize}

\emph{Terms, variables, and substitution.}
Terms are morphisms in the topos. Variables are elements
of objects. Substitution is composition of morphisms. All
are internal.

Therefore: every component of FOL --- connectives,
quantifiers, inference rules, and term structure --- is
a morphism or adjunction within $\mathcal{C}(S,\kappa)$.
The metalanguage is not external to the framework; it is
a fragment of the framework's own internal logic.
\end{proof}

\begin{remark}[The metatheory is internal]
Logicians may object that validating the framework still
requires a metatheory --- rules of inference for reading
and checking the proofs. Proposition \ref{prop:fol} shows
that these rules are themselves derivable within the topos.
The framework's internal logic contains FOL as a fragment,
so the metatheory required to validate the framework is
a sub-theory of the framework's own internal logic. This
does not eliminate the need for a metatheory; it shows
that the needed metatheory is available internally rather
than externally. Combined with self-hosting (Theorem
\ref{thm:self-hosting}) and self-entailment (Theorem
\ref{thm:entailed}), this means: the framework provides
its own ground field, its own primitive operation, and
its own rules of inference.
\end{remark}

\begin{proposition}[Boolean logic as a dependent type]
\label{prop:bool-dependent}
Full Boolean logic over a dimension $d$ is obtained by
imposing the constraint $\chi_d(a) = 1$ for all $a \in S$,
i.e.\ requiring that exactly one of $\tau_d$ or $\sigma_d$
fires on every element. This constraint is itself a
predicate --- a meta-discriminator in the $\sigma$-perspective
--- that classifies discriminators into ``Boolean'' and
``non-Boolean'' based on their behavior over the population.

Since the type of an element's membership profile
(whether it occupies two or four possible states) depends
on a type-level assertion about the discriminator (whether
$\chi_d = 1$ everywhere), this constraint is a
\emph{dependent type} in the system's own type theory.
\end{proposition}

\begin{proof}
The XOR signature $\chi_d(a) = \tau_d(a) + \sigma_d(a)$
is a map $\chi_d: S \to \GF(2)$. The assertion
$\forall a \in S,\; \chi_d(a) = 1$ is a universal
statement about $d$'s behavior on $S$. This statement is
a predicate on $d$ (a discriminator), not on elements of
$S$. Under the $\sigma$-perspective (Proposition
\ref{prop:perspectives}), discriminators are the population,
so this predicate is a $\sigma$-regime discriminator: a
meta-discriminator $\hat{d}_{\mathrm{Bool}}$ with
$\tau_{\hat{d}}(d) = 1$ if $d$ is Boolean over $S$ and
$\sigma_{\hat{d}}(d) = 1$ if $d$ is non-Boolean.

The dependent typing arises because: if
$\hat{d}_{\mathrm{Bool}}$ classifies $d$ as Boolean, then
elements of $S$ under $d$ have membership profiles in
$\{(1,0), (0,1)\}$ (two states). If $d$ is non-Boolean,
profiles range over all of $\GF(2) \times \GF(2)$ (four
states). The type of the membership profile depends on
the type-level fact about the discriminator.
\end{proof}

\begin{remark}[Classical logic is a type, not a restriction]
This result completes the paper's account of Boolean
structure. Classical logic is not:
\begin{itemize}[nosep]
  \item a default that must be restricted (as in ZFC's
    separation axiom),
  \item an alternative logic that competes with
    intuitionistic or paraconsistent logic,
  \item a global property of the system.
\end{itemize}
Classical logic is a \emph{dependent type} within the
system's four-valued logic. It can be freely mixed with
unconstrained reasoning: some dimensions may carry the
Boolean type (enabling classical inference), while others
remain unconstrained (operating in FDE). The type system
tracks which is which, and the boundary between classical
and non-classical reasoning is a first-class object in the
framework --- it is a discriminator in the $\sigma$-regime,
subject to the same algebraic laws as every other
discriminator.

Boolean-constrained and unconstrained reasoning coexist
within a single proof, on a single population, under a
single $\kappa$. The ``choice'' between classical and
non-classical logic is not a foundational commitment; it
is a per-dimension typing judgment.
\end{remark}

\begin{proposition}[Foundational coherence]
\label{prop:coherence}
The framework's epistemological structure is self-reinforcing
without circularity. Two distinct Boolean structures operate
at different levels:
\begin{enumerate}[label=(\roman*),nosep]
  \item \textbf{Base level}: $\GF(2)$ records whether a
    predicate fires ($1$) or is silent ($0$). This is
    definitional --- it is the ground field, the atom of
    observation. It makes no logical claims.
  \item \textbf{Type level}: the Boolean dependent type
    (Proposition \ref{prop:bool-dependent}) asserts that a
    discriminator achieves exhaustive, exclusive partition
    ($\chi_d = 1$ everywhere). This is a structural theorem
    about the discriminator's behavior, constructed from
    base-level observations.
\end{enumerate}
These two levels are mutually reinforcing:
\begin{itemize}[nosep]
  \item The base level makes the type level \emph{expressible}:
    without $\GF(2)$'s ability to record predicate behavior,
    the assertion $\chi_d = 1$ could not be formulated.
  \item The type level makes the base level \emph{interpretable}:
    on dimensions where the Boolean type is satisfied,
    silence ($0$) upgrades from ``no evidence'' to ``false,''
    because exhaustiveness guarantees that absence of
    $\tau$-evidence entails $\sigma$-membership.
  \item The interpretation does not alter the definition:
    $0$ remains silence at the base level. The upgrade to
    ``false'' is additional structure provided by the type
    level, not a redefinition of the ground field.
\end{itemize}
\end{proposition}

\begin{proof}
The non-circularity follows from the asymmetry between the
two levels. The base level ($\GF(2)$ as ground field) exists
independently of any type-level assertion: the exterior
algebra $\bigwedge(\GF(2)^n)$ is defined regardless of
whether any discriminator is Boolean. The type level
(the Boolean dependent type) is constructed from base-level
data: it quantifies over the population's $\chi_d$ values,
which are $\GF(2)$-valued observations.

The feedback from type level to base level is
\emph{interpretive}, not \emph{constitutive}. Knowing that
$d$ is Boolean allows one to \emph{interpret} $\tau_d(a) = 0$
as $\sigma_d(a) = 1$ (by exhaustiveness), but this
interpretation does not change the value of $\tau_d(a)$
in $\GF(2)$. It adds a derived fact ($\sigma_d(a) = 1$)
that was already present in the data but was not
assertible without the type-level guarantee.

Therefore: the base level provides the observations, the
type level provides the interpretation, and the
interpretation validates the observations without
modifying them. This is mutual exhaustive reinforcement,
not circular derivation.
\end{proof}

\begin{remark}[Epistemological bootstrap]
The framework bootstraps its own epistemology. The
ability to observe whether predicates fire ($\GF(2)$) is
used to construct the ability to determine whether the
logic is classical (the Boolean dependent type), which
in turn validates the full interpretive power of the
base-level observations (upgrading silence to falsehood
where appropriate). Each level is necessary for the
other's full functioning, but neither depends on the
other for its \emph{existence}. $\GF(2)$ exists without
the Boolean type; the Boolean type has meaning only
through $\GF(2)$.

This is the same self-supporting structure seen throughout
the paper: sets exist through discrimination
(Section \ref{sec:settheory}), categories exist through
sets (Section \ref{sec:cattheory}), and the logic exists
through the categories (Section \ref{sec:topos}). At no
point does a later construction invalidate or modify an
earlier one. Each layer enriches the interpretation of
the layers below without altering their definitions. The
framework is a tower of mutually reinforcing
interpretations, grounded in $\GF(2)$ and the wedge
product.
\end{remark}

\begin{theorem}[Self-hosting]\label{thm:self-hosting}
The discrimination framework is \emph{self-hosting}:
it can construct its own ground field internally, eliminating
any external dependency on $\GF(2)$ as a given.
\end{theorem}

\begin{proof}
The framework begins with $\GF(2)$ as the ground field of
the exterior algebra (Section \ref{sec:prelim}). From this
ground field, it constructs:
\begin{enumerate}[label=(\roman*),nosep]
  \item the four-valued logic $\GF(2) \times \GF(2)$
    (Definition \ref{def:discriminator}),
  \item the Boolean dependent type, which identifies
    discriminators achieving exhaustive, exclusive
    partition (Proposition \ref{prop:bool-dependent}),
  \item full Boolean logic on dimensions carrying the
    Boolean type (the restriction to $\{(1,0), (0,1)\}$).
\end{enumerate}

Boolean logic over two values is exactly $\GF(2)$: the
field with two elements, addition modulo 2, multiplication
as conjunction. The Boolean dependent type therefore
constructs $\GF(2)$ within the system's own type theory.

The internally constructed $\GF(2)$ has the same algebraic
properties as the externally given $\GF(2)$: it is a field
of characteristic 2, with $1 + 1 = 0$ and $-x = x$.
The exterior algebra $\bigwedge V$ can be defined over
this internal $\GF(2)$ identically to its definition over
the external one.

Therefore the framework can serve as its own substrate.
The construction proceeds:
\[
  \GF(2)_{\mathrm{ext}} \longrightarrow
  \text{four-valued logic} \longrightarrow
  \text{Boolean type} \longrightarrow
  \GF(2)_{\mathrm{int}} \longrightarrow
  \text{four-valued logic} \longrightarrow \cdots
\]
The internal $\GF(2)_{\mathrm{int}}$ is algebraically
identical to $\GF(2)_{\mathrm{ext}}$. The external
dependency can be dropped: the system re-derives its own
ground field from its own type theory at each iteration.
\end{proof}

\begin{remark}[The only primitive is discrimination]
Self-hosting eliminates $\GF(2)$ as an external primitive.
The framework's sole genuine primitive is
\emph{discrimination itself}: the act of applying a
predicate to a population and observing whether it fires.
Everything else is constructed:
\begin{itemize}[nosep]
  \item $\GF(2)$ is the record of a single observation
    (fired or silent) --- and this record is itself
    reconstructible internally via the Boolean type.
  \item The exterior algebra is the structure of composed
    observations.
  \item Sets, categories, types, logic, and the topos are
    the accumulated structure of all observations
    performed.
\end{itemize}
The framework does not assume mathematics and then derive
discrimination. It assumes discrimination and then derives
mathematics --- including, at the end, the ground field
from which it was derived. The self-reference is
non-pathological: $\GF(2)_{\mathrm{int}}$ is
algebraically identical to $\GF(2)_{\mathrm{ext}}$
(Theorem \ref{thm:self-hosting}), so the reconstruction
does not diverge, and $x \wedge x = 0$ prevents the
self-reference from producing a contradictory fixpoint.
\end{remark}

\begin{theorem}[Discrimination is a fixed point]
\label{thm:entailed}
Discrimination is a stable fixed point of the
$\GF(2)$ toroid's closure operation: the 2-cell that,
when applied to the triangle formed by the adjunction,
the ground field, and the product logic, reproduces
itself. No external observer is required. The
contrapositive --- that removing discrimination produces
a structurally unstable gap that reinstates
discrimination --- is proved in Theorem
\ref{thm:irremovable}.

The proof shows discrimination is \emph{a} fixed point;
uniqueness (that no other operation fills the 2-cell)
would require showing the homology class of the gap has
a unique representative, which depends on $H_2 = 0$ for
the toroid's chain complex. We conjecture this holds but
do not prove it here.
\end{theorem}

\begin{proof}
The proof establishes that three structures coexist
coherently and that their binding operation is
discrimination. This is a structural observation about
compatibility --- not a derivation of discrimination from
something prior to it, but a demonstration that
discrimination is the unique operation consistent with
the toroid's geometry.

Consider three structures. Crucially, none is defined in
terms of an observer or an external tool; each is a purely
structural property of a category with projections.

\begin{enumerate}[label=(\roman*),nosep]
  \item \textbf{The adjunction}: the free-forgetful
    adjunction (Theorem \ref{thm:adjunction}) between
    $\kappa$-regimes of different resolution. This is the
    structural capacity to project between finer and
    coarser discrimination regimes. A \emph{predicate} is
    not an external epistemological tool applied by an
    observer; it \emph{is} this adjunction --- a geometric
    projection of the category onto a sub-regime
    (Corollary \ref{cor:free-fgt}). ``Observing whether a
    predicate fires'' is traversing the forgetful functor
    $\pi: \kappa \to \kappa'$: projecting a finer structure
    onto a coarser one and registering whether the
    projected element retains its $\tau/\sigma$ profile.

  \item \textbf{$\GF(2)$}: the algebraic structure of the
    target of the projection. The forgetful functor's
    target is a regime of lower dimension; the minimal
    such target is a single-dimensional regime whose
    $\tau/\sigma$ profile is a single element of
    $\GF(2) \times \GF(2)$, projected to a single
    component: $\GF(2) = \{0,1\}$. $\GF(2)$ is the
    algebraic structure of the simplest possible
    projection.

  \item \textbf{Belnap's four-valued logic}: the product
    $\GF(2) \times \GF(2)$, arising when the full
    (unprojected) $\tau/\sigma$ profile is retained.
    This is the minimal structure that accommodates the
    four possible outcomes of a projection that has not
    been collapsed to a single component: gap $(0,0)$,
    $\tau$-match $(1,0)$, $\sigma$-match $(0,1)$, and
    overlap $(1,1)$.
\end{enumerate}

These three structures are the three non-zero points of
$\GF(2)^2$ (Definition \ref{def:toroid}): the three
vertices of the triangle on the $\GF(2)$ toroid. Each
pair determines the third via the triangle identity
(Definition \ref{def:triangle}):
\begin{itemize}[nosep]
  \item Adjunction $+$ $\GF(2)$ $=$ Belnap
    (a projection given algebraic structure, applied
    independently to two components, yields the
    four-valued product).
  \item $\GF(2)$ $+$ Belnap $=$ Adjunction
    (the algebraic structure of a product, equipped
    with the capacity to project between resolutions,
    yields the adjunction).
  \item Belnap $+$ Adjunction $=$ $\GF(2)$
    (the four-valued logic, projected to a single
    component via the forgetful functor, yields the
    ground field).
\end{itemize}

The three vertices form the boundary of a 2-cell
(Definition \ref{def:boundary}). By the structure of the
discrimination complex, a 2-cell whose boundary consists
of three 1-cells is the wedge product of those 1-cells:
the operation that binds them into a coherent composite.
This binding operation --- the geometric act that mediates
between projection, algebra, and logic --- \emph{is}
discrimination.

No external observer is required at any step. The
adjunction is a property of the category's internal
structure (Theorem \ref{thm:adjunction}). $\GF(2)$ is the
target of the category's own projections. Belnap logic is
the category's own membership profile space. The toroid
``observes itself'' by projecting onto its own sub-spaces
--- this self-projection is the adjunction, and the
adjunction is the predicate.

Therefore: the three structures coexist (each pair entails
the third), none requires an external observer, and the
2-cell that fills their triangle is the unique stable
completion of the toroid's geometry.
That 2-cell is discrimination --- the fixed point of the
toroid's own closure operation.
\end{proof}

\begin{remark}[The predicate is the adjunction]
This formulation eliminates a potential circularity in the
derivation. One might object that ``observing whether a
predicate fires'' presupposes a predicate, which is itself
a discriminator, making the derivation circular. But a
predicate is not an external tool wielded by a homunculus.
It is the free-forgetful adjunction
(Theorem \ref{thm:adjunction}): a structural, geometric
projection native to the $\GF(2)$ toroid. ``Observation''
is the traversal of a functor; ``firing'' is the
retention of a $\tau/\sigma$ profile under projection;
``silence'' is the collapse of that profile. All of these
are internal operations of the category, requiring no
observer external to the system. Formally: the
predicate that the entailment theorem requires is
identical to the adjunction that the entailment theorem
derives. The input and the output of the derivation are
the same mathematical object (the 2-cell filling the
triangle), identified by Theorem \ref{thm:adjunction}.
This is self-reference, but it is non-pathological because
$x \wedge x = 0$: the 2-cell composed with itself is
trivial, preventing the circularity from propagating.
\end{remark}

\begin{remark}[Triadic indivisibility]
Theorem \ref{thm:entailed} does not derive discrimination
from three simpler parts that are \emph{prior} to it.
The adjunction, $\GF(2)$, and Belnap logic are not
building blocks assembled into discrimination; they are
the three irreducible facets of discrimination, which
cannot be separated without destroying the geometry of
the toroid. The triangle identity
$\kappa = \sigma + \tau$ means each facet is the XOR of
the other two: remove one, and the relationship between
the remaining two is undefined.

Discrimination is therefore \emph{mathematically
indivisible} in the precise sense that its internal
triadic structure admits no proper sub-structure that is
self-consistent. The formal content of this claim is not
that the framework was built from nothing, but that
discrimination is the \emph{minimal fixed point} of the
gap-filling operation: it cannot be decomposed into
simpler components that are independently coherent.
\end{remark}

\begin{theorem}[Discrimination is irremovable]
\label{thm:irremovable}
Discrimination cannot be removed from the framework without
the framework deriving it from its own absence. The removal
is self-defeating.
\end{theorem}

\begin{proof}
We prove the contrapositive of Theorem \ref{thm:entailed}:
assume discrimination is absent and derive a contradiction.

\emph{Assume:} the 2-cell binding the three vertices
(adjunction, $\GF(2)$, Belnap logic) does not exist. The
triangle on the $\GF(2)$ toroid has three vertices and
three edges but no filling 2-cell.

\emph{Consequence 1: sheaf failure.}
Without the 2-cell, the covering morphisms between
perspectives (the six adjunctions of Theorem
\ref{thm:adjunction}) do not exist. Local data on
$\kappa$, $\sigma$, $\tau$ cannot be assembled into global
sections: the gluing axiom has no transition maps to
enforce. The site $\mathfrak{S}$ (Definition \ref{def:site})
loses its Grothendieck topology. The category
$\mathcal{C}(S,\kappa)$ ceases to be a sheaf topos and
degrades to three disconnected presheaves --- one per
vertex, with no interaction.

\emph{Consequence 2: manifold failure.}
The $\GF(2)$ toroid is a manifold whose atlas consists of
charts related by the six adjunctions as transition maps.
Without the 2-cell, the transition maps vanish. The three
perspectives become disconnected patches with no atlas.
The toroid is no longer a manifold; it is three isolated
points in $\GF(2)^2$.

\emph{Consequence 3: bimodule failure.}
$\sigma$ serves as a bimodule between $\kappa$ and $\tau$:
it carries a left $\kappa$-action and a right $\tau$-action
simultaneously. The bimodule structure \emph{is} the
discrimination --- it is the operation that allows
$\kappa$-data and $\tau$-data to interact. Without it,
$\kappa$ and $\tau$ are algebraically isolated: they
cannot constrain each other, the triangle identity
$\kappa = \sigma + \tau$ loses its mediating term, and the
relationship between any two vertices is undefined.

\emph{Consequence 4: the gap is self-resolving.}
The unfilled 2-cell is a homological gap: a cycle in
$\ker \partial_1$ that is not in $\operatorname{im}
\partial_2$ (a non-trivial element of $H_1$). By Theorem
\ref{thm:exhaustive} (operational exhaustivity), this gap
fully specifies its own resolution: the precise 2-cell
needed to fill it. That 2-cell has boundary
$\{$adjunction, $\GF(2)$, Belnap$\}$ --- the three edges
of the triangle. Its shape is uniquely determined
(Theorem \ref{thm:exhaustive}, claim (iii), under the
full $\kappa$).

The specified 2-cell \emph{is} discrimination: the
operation that mediates between projection (adjunction),
algebra ($\GF(2)$), and logic (Belnap). Removing
discrimination creates a residue whose self-specified
resolution is discrimination. The removal triggers a
$\kappa$-evolution that reinstates exactly what was
removed.

Therefore: discrimination cannot be absent. Its absence
is a gap; the gap specifies its own filling; the filling
is discrimination. The system derives its own necessity
from its own absence.
\end{proof}

\begin{remark}[Internal necessity]
Theorem \ref{thm:entailed} showed that discrimination
is entailed by the coexistence of its three facets.
Theorem \ref{thm:irremovable} shows the converse: even
if discrimination is hypothetically removed, the
framework's own machinery (homological gap detection,
operational exhaustivity, residue-driven $\kappa$-evolution)
reinstates it. Within the framework, discrimination is
a \emph{stable fixed point}: its removal is its own
reinstatement.

The precise claim is \emph{internal} necessity, not
ontological necessity. The theorems show that the
framework is a fixed point of its own gap-filling
operation, and that this fixed point is unique (the
gap's shape specifies the filling). This is stronger
than mere consistency (Peano arithmetic is consistent but
not self-hosting) and stronger than mere self-description
(PA can describe its own signature but cannot reconstruct
its axioms from a gap). The framework reconstructs its
own primitive operation from the structural instability
of its absence.

Whether ``internally necessary'' and ``ontologically
necessary'' are themselves distinguishable is a question
the framework addresses on its own terms. No
discriminator within the framework separates the two
claims: any test for ``ontological necessity'' would be a
discrimination, which is the operation whose necessity is
at issue. By $\kappa$-extensionality, if no discriminator
separates the claims, they are identical \emph{within the
framework}. The framework dissolves the objection rather
than deflecting it --- but this dissolution is available
only to readers who have accepted the operationalist
principle (Remark \ref{rem:foundational-status}), which
is the one element the framework explicitly flags as a
foundational choice rather than a derivation.
\end{remark}

\begin{remark}[The adjunction as formal mechanism]
The triadic adjunction system (Theorem \ref{thm:adjunction})
provides the formal mechanism for this entailment. The six
permutations of $\{\kappa, \sigma, \tau\}$ over the
adjunction stencil $A \leftarrow B \to C$ are precisely
the six ways to assign the roles of left adjoint, hom, and
right adjoint to the three facets. The triangle identity
guarantees that each assignment produces a valid adjunction
--- that is, that any two of the three determine the third.

The 2-cell that fills the triangle is the adjunction itself:
the natural bijection $\operatorname{Hom}_B(A, C) \cong
A + C$ that the triangle identity provides. This bijection
is the operation that mediates between the three structures.
That operation is discrimination. The triadic adjunction
does not merely use the triangle identity; it \emph{is}
the triangle identity, applied to the foundational triangle.

Chirality determines interpretation. The two adjunctions on
each axis --- $\kappa \leftarrow \sigma \to \tau$ and
$\tau \leftarrow \sigma \to \kappa$ --- correspond to two
orientations of the same mediation. The dual to a functor
(reversing chirality) reverses which structure is ``input''
and which is ``output'' without changing the underlying
relationship. The triangle identity
$\kappa = \sigma + \tau$ ensures that all three
perspectives access the same information (each is
determined by the other two, preserving dimension).
What changes under rotation is not the content but the
perspective from which the content is accessed.
\end{remark}

\begin{remark}[Internal recoverability]
Theorem \ref{thm:self-hosting} showed that the framework
reconstructs its own ground field internally.
Theorem \ref{thm:entailed} shows that it reconstructs
its own primitive operation. The framework's external
dependencies are \emph{recoverable from within}:
\begin{itemize}[nosep]
  \item Discrimination is a stable fixed point of the
    framework's own gap-filling operation
    (Theorem \ref{thm:irremovable}).
  \item $\GF(2)$ is reconstructible internally via the
    Boolean dependent type.
  \item The exterior algebra is the structure of composed
    discriminations.
  \item Sets, categories, types, logic, and the topos are
    accumulated structure.
  \item The four-cell structure that began as a definition
    turns out to be the subobject classifier.
  \item The toroid that organizes the perspectives turns
    out to be the space whose closure entails
    discrimination.
\end{itemize}
The framework is \emph{internally closed}: every
component it requires is derivable from within. This is
not the claim that the framework ``assumes nothing'' ---
the operationalist axiom (Definition
\ref{def:operationalist}) is a foundational
commitment, and the initial articulation of any predicate
is an act that the framework describes but does not
conjure from the void. The claim is that the framework
is a \emph{fixed point of its own operations}: once any
discrimination is performed, the entire structure is
recoverable, including the ground field and the
discrimination operation itself.
\end{remark}

\subsection{Sheaf-theoretic interpretation}\label{sec:sheaf}

The triadic adjunction system (Theorem \ref{thm:adjunction})
provides the sheaf-theoretic structure of
$\mathcal{C}(S,\kappa)$ by witnessing the construction of
every dimension of the toroid.

\begin{definition}[The site of regimes]\label{def:site}
The \emph{site of regimes} $\mathfrak{S}$ is the category
whose:
\begin{itemize}[nosep]
  \item objects are $\kappa$-regimes of all dimensions
    (sub-regimes $\kappa' \subset \kappa$ obtained by
    projection),
  \item morphisms are projections $\pi: \kappa \to \kappa'$
    (forgetful functors) and inclusions
    $\iota: \kappa' \hookrightarrow \kappa$ (free functors).
\end{itemize}
The \emph{Grothendieck topology} on $\mathfrak{S}$ is:
a family of projections $\{\pi_i: \kappa \to \kappa'_i\}$
\emph{covers} $\kappa$ if the union of the sub-regimes
$\kappa'_i$ recovers the full $\tau/\sigma$ information
of $\kappa$ --- that is, if every discriminator in $\kappa$
belongs to at least one $\kappa'_i$.
\end{definition}

\begin{theorem}[The topos is a sheaf topos]\label{thm:sheaf}
$\mathcal{C}(S,\kappa)$ is equivalent to the category of
sheaves on the site of regimes $\mathfrak{S}$.
\end{theorem}

\begin{proof}
A \emph{presheaf} on $\mathfrak{S}$ assigns to each regime
$\kappa'$ a collection (the set of $\kappa'$-equivalence
classes of $S$) and to each projection $\pi: \kappa \to
\kappa'$ a restriction map (coarsening the equivalence
classes). A presheaf is a \emph{sheaf} if it satisfies the
gluing axiom: compatible data on overlapping sub-regimes
glue uniquely to data on the union.

The gluing axiom is the order-independence of filtration
(Proposition \ref{prop:monoidal}). If two sub-regimes
$\kappa'_1$ and $\kappa'_2$ produce $\tau/\sigma$
assignments that agree on their shared discriminators, the
monoidal property guarantees that the joint filtration
through $\kappa'_1 \cup \kappa'_2$ is well-defined and
unique. This is precisely the sheaf condition: compatible
local data (partitions on sub-regimes) glue to global data
(a partition on the union).

Every object of $\mathcal{C}(S,\kappa)$ defines a sheaf
on $\mathfrak{S}$: an element $a \in S$ assigns to each
sub-regime $\kappa'$ its $\tau/\sigma$ profile under
$\kappa'$, and the restriction maps are the projections'
action on profiles. Conversely, every sheaf on $\mathfrak{S}$
that satisfies the representability condition
(compatible local profiles for some element) corresponds
to an object of $\mathcal{C}(S,\kappa)$.

The equivalence of categories follows: the Yoneda
embedding (Theorem \ref{thm:yoneda}) embeds
$\mathcal{C}(S,\kappa)$ fully and faithfully into
presheaves on $\mathfrak{S}$, and the sheaf condition
(Proposition \ref{prop:monoidal}) selects exactly the
representable presheaves.
\end{proof}

\begin{remark}[The six adjunctions are the covering data]
Each of the six triadic adjunctions
(Theorem \ref{thm:adjunction}) is a covering morphism in
$\mathfrak{S}$. The three axes of the toroid correspond to
three independent ``directions'' of projection, and the two
chiralities on each axis correspond to the two orientations
of each covering. The six adjunctions jointly provide the
\emph{complete} covering data for the site: every
projection between regimes factors through a composition
of these six fundamental covering morphisms.

This means the Grothendieck topology on $\mathfrak{S}$ is
generated by the triadic adjunction system. The site is
not an external structure imposed on the category; it is
the toroid's own geometry, expressed as covering data.
\end{remark}

\begin{remark}[The sheaf condition was already proved]
The identification of $\mathcal{C}(S,\kappa)$ as a sheaf
topos is not a new result. It is the recognition that
Proposition \ref{prop:monoidal} (order-independence of
filtration) \emph{is} the sheaf condition, and that the
site $\mathfrak{S}$ \emph{is} the category of
$\kappa$-regimes with the projection topology. The sheaf
structure was present in the framework from Section 3;
it merely required the toroid geometry of Sections 7--8
to be recognized as a Grothendieck topology.
\end{remark}

\begin{proposition}[The four-cell structure is the
  presheaf-to-sheaf transition]\label{prop:presheaf-sheaf}
The four cells of the membership profile
(Definition \ref{def:discriminator}) correspond exactly
to the stages of the sheaf gluing process:
\begin{center}
\begin{tabular}{c|l|l}
$(\tau, \sigma)$ & Sheaf status & Covering status \\
\hline
$(0,0)$ & Presheaf & No covering morphism interrogated \\
$(1,0)$ & Sheaf & $\tau$-covering sufficient \\
$(0,1)$ & Sheaf & $\sigma$-covering sufficient \\
$(1,1)$ & Over-covered & Coverings conflict; residue
\end{tabular}
\end{center}
A value of $0$ on a component means the corresponding
adjunctive perspective has not been interrogated on this
dimension --- the covering morphism has not been applied.
A value of $1$ means data has been obtained from that
perspective.
\end{proposition}

\begin{proof}
A presheaf on $\mathfrak{S}$ assigns data to sub-regimes
without requiring the gluing condition. The gap state
$(0,0)$ is the maximally presheaf condition: no projection
has produced data on this dimension. The element exists in
the population but is invisible to every covering morphism
that has been applied.

The Boolean states $(1,0)$ and $(0,1)$ are sheaf states:
a sufficient covering has been applied (the $\tau$- or
$\sigma$-adjunction has been interrogated), the gluing
axiom is satisfied (there is a unique, consistent
assignment), and the element's profile on this dimension
is determined.

The overlap state $(1,1)$ is the over-covered state: two
independent coverings have both produced data, and the
data is conflicting (the element matches both $\tau$ and
$\sigma$). The gluing axiom fails at this point --- the
local data does not assemble into a unique global section.
This failure is a residue (Definition \ref{def:residue}),
which triggers $\kappa$-evolution: a refinement of the
topology that resolves the conflict by introducing a
finer covering.

The transition from presheaf to sheaf is therefore the
transition from $(0,0)$ to $(1,0)$ or $(0,1)$: the
interrogation of covering morphisms until the data is
sufficient and consistent. $\kappa$-evolution drives this
transition by providing new covering morphisms (new
adjunctive perspectives) that resolve gaps and conflicts.
\end{proof}

\begin{remark}[Convergence as sheafification]
The dynamic proof theory's tendency toward Boolean
(Theorem \ref{thm:proof-theory}) is \emph{sheafification}:
the process of converting presheaves into sheaves by
applying covering morphisms until the gluing axiom is
satisfied. Each $\kappa$-evolution interrogates a new
adjunctive perspective, moving elements from the presheaf
state $(0,0)$ or the over-covered state $(1,1)$ toward
the sheaf states $(1,0)$ and $(0,1)$.

Sheafification is the left adjoint to the inclusion of
sheaves into presheaves --- and this adjunction is itself
an instance of the triadic adjunction system
(Theorem \ref{thm:adjunction}). The framework's dynamics
are sheafification, and sheafification is an adjunction,
and the adjunction is a rotation of the toroid.

The two remaining states that resist sheafification ---
elements stuck at $(0,0)$ because no covering reaches
them, or at $(1,1)$ because every covering over-covers
them --- are precisely the residues that drive
$\kappa$-evolution. Sheafification never terminates for
an inexhaustible collection (Proposition
\ref{prop:infinity}), but every individual dimension
can be sheafified in finite time.
\end{remark}

\subsection{The expanded classifier: $\GF(2)^3$}
\label{sec:expanded}

The presheaf-to-sheaf identification (Proposition
\ref{prop:presheaf-sheaf}) reveals an ambiguity in the
original four-cell structure: the gap state $(0,0)$
conflates two distinct situations:
\begin{itemize}[nosep]
  \item ``The perspective was interrogated and produced no
    match'' (genuine silence --- data was sought and not
    found).
  \item ``The perspective has not been interrogated at all''
    (no covering morphism has been applied --- no data was
    sought).
\end{itemize}
Distinguishing these requires a third dimension.

\begin{definition}[Coverage profile]\label{def:coverage}
For each of the three perspectives $\kappa$, $\sigma$,
$\tau$, define an independent Boolean:
\[
  \gamma_\kappa, \gamma_\sigma, \gamma_\tau \in \GF(2),
\]
where $\gamma_x = 1$ means ``the $x$-perspective has been
interrogated (the corresponding covering morphism has been
applied and there is something to report)'' and
$\gamma_x = 0$ means ``the $x$-perspective has not been
interrogated (no covering has been applied from this
perspective).''

The \emph{full membership profile} of an element $a$
under discriminator $d$ in regime $\kappa$ is a triple
\begin{equation}
  (\gamma, \tau_d(a), \sigma_d(a)) \in \GF(2)^3,
\end{equation}
where $\gamma$ encodes the coverage status and
$(\tau_d(a), \sigma_d(a))$ encodes the content.
The content is meaningful only when $\gamma = 1$;
when $\gamma = 0$, the content bits are undefined.
\end{definition}

\begin{proposition}[The three-dimensional classifier]
\label{prop:gf2-cubed}
The expanded subobject classifier is
$\Omega' = \GF(2)^3$, with five operationally distinct
states:
\begin{center}
\begin{tabular}{c|c|l|l}
$\gamma$ & $(\tau, \sigma)$ & State & Sheaf status \\
\hline
$0$ & --- & Uncovered & Pure presheaf \\
$1$ & $(0,0)$ & Covered, silent & Interrogated gap \\
$1$ & $(1,0)$ & $\tau$-match & Sheaf (determined) \\
$1$ & $(0,1)$ & $\sigma$-match & Sheaf (determined) \\
$1$ & $(1,1)$ & Both match & Over-covered (residue)
\end{tabular}
\end{center}
The ``uncovered'' state $(\gamma = 0)$ is the true
presheaf condition: the covering morphism has not been
applied. The ``covered, silent'' state
$(\gamma = 1, \tau = \sigma = 0)$ is a genuinely
interrogated gap: the covering was applied and the
element was invisible to it. These are operationally
distinct and must not be conflated.
\end{proposition}

\begin{proof}
The distinction is forced by the sheaf interpretation.
A presheaf assigns data to open sets; a section is
\emph{defined} on an open set $U$ if $U$ has been
covered, and \emph{undefined} if $U$ has not. The
sheaf condition applies only to defined sections.
If $\gamma = 0$, no section exists to glue; the element
is outside the domain of the covering. If $\gamma = 1$
and $(\tau, \sigma) = (0,0)$, a section exists but
carries no data --- the covering was applied but
produced no match.

The original $\GF(2)^2$ classifier (Definition
\ref{def:subobj-class}) is the $\gamma = 1$ slice of
$\GF(2)^3$. All results proved using $\Omega = \GF(2)^2$
remain valid within this slice. The expansion adds the
$\gamma = 0$ state without modifying the existing
structure.
\end{proof}

\begin{proposition}[Gluing mechanics across coverage states]
\label{prop:gluing-mechanics}
The sheaf gluing axiom operates on the $\GF(2)^3$
classifier as follows. Let $d$ be a discriminator and
$a$ an element with coverage-content profile
$(\gamma, \tau, \sigma)$.

\emph{Step 1: Coverage interrogation.}
Before any covering morphism is applied, $\gamma = 0$.
The content $(\tau, \sigma)$ is undefined. The element
$a$ on dimension $d$ is in the pure presheaf state:
it exists in the population but has not been probed.

\emph{Step 2: Applying a covering morphism.}
The forgetful functor $\pi: \kappa \to \kappa'$
(one of the six triadic adjunctions) is applied. This
sets $\gamma = 1$ and evaluates the content: the
discriminator $d$ produces $(\tau_d(a), \sigma_d(a))$.
The result is one of the four Belnap states:
$(1,0)$, $(0,1)$, $(0,0)$, or $(1,1)$.

\emph{Step 3: Gluing condition.}
Suppose two covering morphisms $\pi_1, \pi_2$ (from
two different triadic adjunctions) are applied to $a$
on overlapping sub-regimes. Each produces a
coverage-content triple:
$(\gamma_1 = 1, \tau_1, \sigma_1)$ and
$(\gamma_2 = 1, \tau_2, \sigma_2)$.
The gluing axiom requires: on the overlap (the
discriminators common to both sub-regimes),
$(\tau_1, \sigma_1) = (\tau_2, \sigma_2)$. This holds
by Proposition \ref{prop:monoidal} (order-independence):
the $\tau/\sigma$ profile is a function of $a$ and $d$
alone, independent of which covering morphism was applied.

\emph{Step 4: Fano inference on $\gamma = 0$ components.}
If two linearly independent perspectives have $\gamma = 1$
and a third has $\gamma = 0$, the Fano line through the
two interrogated points determines the third
(Corollary \ref{cor:self-model}). The inferred content is
$p_3 = p_1 + p_2$ (the triangle identity applied
componentwise). The third component transitions from
$\gamma = 0$ (uncovered) to $\gamma = 1$ (covered by
inference) without direct interrogation.

\emph{Step 5: Sheaf completion.}
After three linearly independent interrogations (Theorem
\ref{thm:convergence}), all seven Fano points have
$\gamma = 1$. The sheaf is fully glued: every component
has well-defined content, and the gluing axiom is
satisfied by order-independence. The element's full
profile on dimension $d$ is determined.
\end{proposition}

\begin{proof}
Steps 1--3 follow from Definition \ref{def:coverage} and
Proposition \ref{prop:monoidal}. Step 4 follows from
Theorem \ref{thm:fano} and Corollary \ref{cor:self-model}.
Step 5 follows from Theorem \ref{thm:convergence}.

The key verification is Step 4: is the inferred content
consistent with what direct interrogation would produce?
Yes: the triangle identity $p_3 = p_1 + p_2$ holds
because $\kappa = \sigma + \tau$ (Definition
\ref{def:triangle}), and this identity governs the
relationship between the three perspectives at every
level (the content level, the coverage level, and the
adjunction level). Inference via Fano lines computes
the same value that direct interrogation would compute,
because both are governed by the same algebraic law.
\end{proof}

\begin{remark}[Coverage as dependent type]
The coverage bit $\gamma$ is itself a dependent type
in the same sense as the Boolean type (Proposition
\ref{prop:bool-dependent}): it is a predicate on the
discriminator (``has this discriminator been applied to
this element?''), not on the element itself. The type
of the content $(\tau, \sigma)$ depends on $\gamma$:
if $\gamma = 1$, the content is a well-defined element
of $\GF(2)^2$; if $\gamma = 0$, the content is
undefined.

This is the framework's account of
\emph{definedness} --- the distinction between ``no
value'' and ``the value zero.'' In classical
mathematics, every variable has a value. In the
discrimination framework, a variable has a value only
on dimensions where the covering has been applied.
Definedness is not a metalinguistic property; it is
a first-class dimension of the membership profile,
tracked by $\gamma$ and governed by the same $\GF(2)$
algebra as every other dimension.
\end{remark}

\begin{theorem}[Fano closure via the triangle identity]
\label{thm:fano}
The seven non-zero points of $\GF(2)^3$ form the Fano
plane $PG(2, \GF(2))$: the projective plane over $\GF(2)$,
with seven points and seven lines, each line passing through
three points. Each line of the Fano plane is an instance
of the triangle identity: three elements that pairwise-XOR
to produce the third. The Fano structure is the complete
set of triangle identities in $\GF(2)^3$.
\end{theorem}

\begin{proof}
In $\GF(2)^3$, three non-zero points $p_1, p_2, p_3$ are
collinear (lie on a line of $PG(2, \GF(2))$) if and only
if $p_1 + p_2 + p_3 = 0$, equivalently
$p_3 = p_1 + p_2$. This is the triangle identity: any
two determine the third. The seven lines of the Fano
plane are therefore seven triangle identities.

The three ``fundamental'' triangle identities of the
framework --- $\kappa = \sigma + \tau$ and its two
permutations --- correspond to three of the seven lines
(those within the $\gamma = 1$ slice, where coverage is
established). The remaining four lines involve mixed
coverage-and-content states, providing triangle identities
that relate coverage status to content.
\end{proof}

\begin{corollary}[Self-modeling]\label{cor:self-model}
If two perspectives lying on the same Fano line have been
interrogated ($\gamma = 1$), the third perspective's
response is determined by the triangle identity on that
line, even if the third perspective has not been directly
interrogated ($\gamma = 0$).

Therefore, the $\gamma = 0$ state is a genuine presheaf
state \emph{only when no Fano line connects it to two
interrogated perspectives}. The unique point where no
inference is possible is $(0,0,0)$: the zero element of
$\GF(2)^3$, the only point \emph{not} on the Fano plane.
This is the sole true presheaf state --- the state of
absolute non-interrogation.

Every other state is either directly interrogated or
inferrable from two interrogated perspectives via a
Fano line. The system models its own uninterrogated
states using the triangle identity applied to its own
coverage data.
\end{corollary}

\begin{proof}
Let $p_1, p_2$ be two interrogated points on a Fano line
$\ell$. The third point on $\ell$ is $p_3 = p_1 + p_2$.
The value at $p_3$ is determined by the values at $p_1$
and $p_2$ via the triangle identity, regardless of whether
$p_3$ has been directly interrogated.

The zero element $(0,0,0)$ does not lie on any line of
$PG(2, \GF(2))$ (projective lines consist of non-zero
points). Therefore no pair of interrogated perspectives
can reach $(0,0,0)$ by inference. It is the unique
irreducible presheaf state.
\end{proof}

\begin{remark}[Self-modeling is the triangle identity at
  the coverage level]
The self-modeling property (Corollary \ref{cor:self-model})
is not a new principle. It is the triangle identity ---
the same identity that governs the content level
($\kappa = \sigma + \tau$), the adjunction structure
(Theorem \ref{thm:adjunction}), the monoidal closure
(Theorem \ref{thm:closure}), and the entailment of
discrimination (Theorem \ref{thm:entailed}) --- now
operating at the coverage level. The system predicts its
own unobserved states using the same algebraic law that
governs everything else.

This is the deepest expression of self-hosting
(Theorem \ref{thm:self-hosting}): the framework not only
constructs its own ground field and its own primitive
operation, but \emph{infers its own uninterrogated
coverage} from its interrogated coverage. The system
models itself, using itself, within itself. The Fano
plane is the geometry of this self-modeling: its seven
lines are the seven inference channels through which
the framework fills its own gaps.
\end{remark}

\begin{theorem}[Convergence rate]\label{thm:convergence}
Sheafification of a single $\kappa$-dimension converges in
exactly $3$ interrogations: the dimension of $\GF(2)^3$.
Once three linearly independent perspectives have been
interrogated on a dimension, all seven Fano points are
determined and the sheaf on that dimension is fully glued.
\end{theorem}

\begin{proof}
A basis of $\GF(2)^3$ consists of $3$ linearly independent
vectors. Once three such vectors are interrogated (their
$\gamma = 1$ and their $(\tau, \sigma)$ content is known),
every other element of $\GF(2)^3$ is a linear combination
of these three. Since the Fano lines are the triangle
identities $p_3 = p_1 + p_2$, and every non-zero point is
reachable from any basis by at most one addition, the
remaining four points are determined.

Therefore: three linearly independent interrogations per
dimension suffice for complete coverage. The convergence
rate is $3 / 7$ (three interrogations determine all seven
perspectives), and this is sharp: fewer than three linearly
independent interrogations leave at least one perspective
undetermined.
\end{proof}

\begin{corollary}[The convergence rate is the triangle]
\label{cor:convergence-triangle}
The number $3$ --- the convergence rate, the dimension of
the classifier, the number of vertices of the triangle,
the number of fundamental adjunctions per chirality, and
the number of perspectives in Proposition
\ref{prop:perspectives} --- is the same $3$ throughout the
framework. Convergence is not an asymptotic property; it
is a structural constant determined by the triangle
identity.
\end{corollary}

